% Part: computability
% Chapter: recursive-functions
% Section: normal-form

\documentclass[../../../include/open-logic-section]{subfiles}

\begin{document}

% SEGMENT OLP-0225-S01

\olfileid{cmp}{rec}{nft}
\olsection{இயல்புவடிவத் தேற்றம்}

\begin{thm}[க்ளீனியின் இயல்புவடிவத் தேற்றம்]
\ollabel{thm:kleene-nf}
பின்வரும் பண்புடைய முதன்மை மீள்வரையறை உறவு $T(e, x, s)$ ஒன்றும்
முதன்மை மீள்வரையறைச் சார்பு $U(s)$ ஒன்றும் உள்ளன: $f$ ஏதேனும்
பகுதி மீள்வரையறைச் சார்பு எனில், ஏதேனும்~$e$ க்கு, ஒவ்வொரு
$x$ க்கும்
\[
f(x) \simeq U(\umin{s}{T(e, x, s)})
\]
பொருந்தும்.
\end{thm}

\begin{explain}
இயல்புவடிவத் தேற்றத்தின் நிரூபிப்பு சிக்கலானது; ஆனால் அடிப்படைக்
கருத்து எளியது. ஒவ்வொரு பகுதி மீள்வரையறைச் சார்புக்கும் ஓர்
\emph{சுட்டெண்}~$e$ உண்டு; உள்ளுணர்வாக, அது அச்சார்பின் நிரலை
அல்லது வரையறையைக் குறியீடாக்கும் எண். $f(x) \fdefined$ எனில்,
கணக்கீட்டை முறையாகப் பதிவுசெய்து ஏதேனும் எண்~$s$ ஆல் குறியீடாக்கலாம்;
மேலும் $s$ ஆனது உள்ளீடு~$x$ இல்~$f$ இன் கணக்கீட்டைக் குறியீடாக்குகிறது
என்ற உண்மையை $x$ மற்றும் வரையறை~$e$ ஆகியவற்றை மட்டும் பயன்படுத்தி
முதன்மை மீள்வரையறை முறையில் சோதிக்கலாம். இதன் விளைவாக, உறவு~$T$,
“சுட்டெண்~$e$ உடைய சார்புக்கு உள்ளீடு~$x$ இல் ஒரு கணக்கீடு உள்ளது;
மேலும் $s$ இக்கணக்கீட்டைக் குறியீடாக்குகிறது,” என்பது முதன்மை
மீள்வரையறை உறவு. கணக்கீட்டின் முழுப் பதிவு~$s$ தரப்பட்டால்,
~$s$ இன் “விளைவு”~$f(x)$ இன் மதிப்பு; அதையும்~$s$ இலிருந்து
முதன்மை மீள்வரையறை முறையில் பெறலாம்.

எந்த பகுதி மீள்வரையறைச் சார்பின் வரையறைக்கும் ஒரே ஒரு வரம்பற்ற
தேடல் மட்டும் தேவை என்பதை இயல்புவடிவத் தேற்றம் காட்டுகிறது.
அடிப்படையில், சுட்டெண்~$e$ உடைய சார்பின் உள்ளீடு~$x$ க்கான
கணக்கீட்டைக் குறியீடாக்கும் எண்ணைக் காணும் வரை எல்லா எண்களிலும்
தேடலாம். எண்கள்~$e$ ஐப் பகுதி மீள்வரையறைச் சார்புகளின்
“பெயர்களாகப்” பயன்படுத்தி, தேற்றத்திலுள்ள சமன்பாட்டால் வரையறுக்கப்படும்
சார்பு~$f$ க்கு $\cfind{e}$ என எழுதலாம். எந்த பகுதி மீள்வரையறைச்
சார்பும் ஒன்றுக்கு மேற்பட்ட சுட்டெண்களைக் கொண்டிருக்கலாம் என்பதைக்
கவனிக்க---உண்மையில் ஒவ்வொரு பகுதி மீள்வரையறைச் சார்பும் முடிவிலி
எண்ணிக்கையிலான சுட்டெண்களைக் கொண்டுள்ளது.
\end{explain}
\end{document}
